(1) 取 \(\boldsymbol m=m(\sin\alpha\cos\Omega t,\sin\alpha\sin\Omega t,\cos\alpha)\)。其二阶导数仅横向，\(|\ddot{\boldsymbol m}|^2=m^2\Omega^4\sin^2\alpha\)，故
\[
P_{\rm rad}=K\Omega^4,\qquad K=\frac{\mu_0m^2\sin^2\alpha}{6\pi c^3}.
\]
轴向磁矩 \(\alpha=0\) 不随时间改变，理想偶极辐射为零。
(2) \(E_{\rm rot}=I\Omega^2/2\)，能量平衡 \(I\Omega\dot\Omega=-K\Omega^4\)，即 \(\dot\Omega=-(K/I)\Omega^3\)。分离变量得
\[
\Omega(t)=\frac{\Omega_0}{\sqrt{1+2K\Omega_0^2t/I}}.
\]
慢变近似要求每个自转周期的相对变化 \(|\dot\Omega|2\pi/\Omega^2=2\pi K\Omega/I\ll1\)。
(3) 再微分 \(\ddot\Omega=3(K/I)^2\Omega^5\)，故 \(n=3\)。由 \(P_{\rm rot}=2\pi/\Omega\) 得
\[
\dot P_{\rm rot}=\frac{2\pi K}{I}\Omega,\qquad P_{\rm rot}\dot P_{\rm rot}=\frac{4\pi^2K}{I}.
\]
若 \(I,m,\alpha\) 随时间变、存在粒子风或磁层电流，实测制动指数可不等于3。
(4) \(t_{0\to1}=\frac I{2K}(\Omega_1^{-2}-\Omega_0^{-2})\)，辐射能量 \(\int P_{\rm rad}dt=I(\Omega_0^2-\Omega_1^2)/2\)，等于转动能减少。辐射角动量流对应制动力矩 \(N=P_{\rm rad}/\Omega=K\Omega^3\)，量纲为 \(J\)。本模型采用真空远区磁偶极近似，不能代替等离子体填充磁层的完整自洽制动理论。
